\section{The reachability side: a measured negative}
\label{sec:deferral}

This section reports a failure. It is kept in full because it closes an axis the project has spent
effort on twice, and because the way it failed is more general than the theorem that predicted it.

\subsection{The tax, and the obstacle}

Prior work establishes that $\texttt{explored} - \texttt{stored}$ equals, to the state, the number
of \textbf{cross-level retirements}: a zone stored at level $d$, expanded, then superseded at level
$d+g$ by a more general zone for the same discrete state. The first expansion was wasted. That is
$41.7\%$ of all expansions on \texttt{fischer10} and $40.8\%$ on \texttt{twolock10}, and exactly
zero on \texttt{arbiter}, \texttt{viking} and \texttt{csma}. \measured{}

It also establishes the obstacle. Under a level-synchronous schedule that inserts a depth-$d$ state
at $T(d)$ and expands it at $T(d)+W$, we get $T(d) = d(W+1)$, so the retiring zone at depth $d+g$
arrives at $T(d) + g(W+1)$ and is caught only if
\begin{equation}
  g(W+1) < W, \label{eq:nodelay}
\end{equation}
false for every $g \ge 1$, $W \ge 0$. \proved{} The prior study concluded that only an
\emph{oracle} naming the doomed states could collect the tax, and left one question undecided as
deciding the whole idea.

\subsection{The undecided question, answered: causal independence holds}

Is the retiring zone causally independent of the retired one? If the bigger zone at $q$ arrives
along a path through $q$'s own expansion, parking $q$ parks its own retirer.

We instrumented a BFS with $aLU$ subsumption so every node records its BFS parent and birth level,
and on each retirement walked the parent chain of the state producing the enlarging successor.

\begin{center}\small
\begin{tabular}{lrrrrl}
\toprule
model & stored & expansions & cross-level retirements & \textbf{independent} & gap histogram \\
\midrule
fischer4 & 222    & 270    & 48    & \textbf{48 (100\,\%)}   & 0:12\ \ 4:48 \\
fischer5 & 762    & 1\,012 & 250   & \textbf{250 (100\,\%)}  & 0:110\ \ 4:250 \\
fischer6 & 2\,694 & 3\,774 & 1\,080& \textbf{1080 (100\,\%)} & 0:630\ \ 4:1080 \\
fischer7 & 9\,798 & 14\,012& 4\,214& \textbf{4214 (100\,\%)} & 0:2912\ \ 4:4214 \\
twolock5 & 922    & 1\,172 & 250   & \textbf{250 (100\,\%)}  & 0:110\ \ 6:250 \\
twolock7 & 10\,694& 14\,908& 4\,214& \textbf{4214 (100\,\%)} & 0:2912\ \ 6:4214 \\
phil6/8/10 & --- & --- & \textbf{0} & --- & 0 only \\
arbiter5 & 112 & 112 & \textbf{0} & --- & no retirements at all \\
\bottomrule
\end{tabular}
\end{center}
\measured{}

The answer is favourable: $100\%$ independent at every size, with the gap a model constant ($4$ on
\texttt{fischer}, $6$ on \texttt{twolock}) exactly reproducing the prior study's signature. So the
blocking objection --- ``parking a state stalls its own retirer'' --- is \textbf{false} on these
families.

\subsection{There is no cheap oracle}

Per-\texttt{(process, location)} retirement rates on \texttt{fischer5}: \texttt{Pi.idle} $0.217$,
\texttt{Pi.req} $0.394$, \texttt{Pi.wait} $0.380$, \texttt{Pi.cs} $0.000$, identical across all five
symmetric processes. $32.8\%$ of stored nodes are retired and \textbf{no location-vector feature
separates them}. \measured{} A predictor that had to inspect zones would cost what it saves.

\subsection{The reformulation, and why it fails}

Since prediction is unavailable, the natural move is to \emph{observe} instead: park a popped node,
let other branches run, and if its zone grows while parked, drop the parked entry --- the wasted
expansion is never performed. Inequality~(\ref{eq:nodelay}) rules this out under a level-synchronous
schedule, but it appears not to rule it out on a \emph{barrier-free} frontier, and the independence
result above appears to license exactly that.

We implemented it. The frontier is a single global FIFO with no level barrier --- an asynchronous
frontier observed in one interleaving. Queue entries are (node, version-at-push, park-count); a pop
whose recorded version is stale is dropped (the saving), a pop with park-count $< D$ is re-queued
with the count incremented (the parking), otherwise the node is expanded.

\begin{center}\small
\begin{tabular}{lrrrrrrr}
\toprule
model & level-sync & $D{=}0$ & $D{=}1$ & $D{=}2$ & $D{=}4$ & $D{=}8$ & $D{=}16$ \\
\midrule
fischer5 & 1\,012 & 990    & 990    & 990    & 990    & 990    & 990 \\
fischer6 & 3\,774 & 3\,514 & 3\,514 & 3\,514 & 3\,514 & 3\,514 & 3\,514 \\
twolock5 & 1\,172 & 1\,150 & 1\,150 & 1\,150 & 1\,150 & 1\,150 & 1\,150 \\
twolock6 & 4\,158 & 3\,898 & 3\,898 & 3\,898 & 3\,898 & 3\,898 & 3\,898 \\
\bottomrule
\end{tabular}
\end{center}
\measured{} (expansions; $D$ is the parking budget)

\textbf{Deferral saves exactly zero, at every budget, on every model.} The columns are not close;
they are identical. Parking raised the number of re-queue operations from $0$ to $56\,224$ on
\texttt{fischer6} and changed the expansion count by nothing at all.

\begin{proposition}[The no-op is more general than (\ref{eq:nodelay}) suggests]
\label{prop:uniform}
A \emph{uniform} parking policy is a no-op under \textbf{any} schedule, not merely a
level-synchronous one. Parking every popped node by the same budget $D$ applies the same shift to
the retiree and to its retirer, so their relative order in the queue is unchanged and no
supersession that was not already observed becomes observable. \proved{}, and \measured{} above on a
barrier-free FIFO.
\end{proposition}

This is where the reformulation dies. ``Observe instead of predict'' silently reintroduced
uniformity: waiting a fixed amount for \emph{every} node is a uniform delay by another name.
Removing the barrier does not help, because the barrier was never what made
(\ref{eq:nodelay}) true --- uniformity was.

\subsection{Status of this axis}

Both routes are now closed by measurement rather than by argument:

\begin{itemize}
  \item \textbf{Uniform deferral} --- no-op, Proposition~\ref{prop:uniform}, measured at six
        budgets on four models.
  \item \textbf{Selective deferral} --- requires a per-state oracle; no location-vector predictor
        exists, and a zone-level one costs what it saves.
\end{itemize}

The causal-independence result stands and is worth recording, because it removes the objection that
was thought to be fatal --- but removing that objection turns out not to be sufficient. The one
genuinely encouraging number is incidental: the barrier-free FIFO expands $2.2\%$--$6.9\%$ fewer
states than level-synchronous BFS ($1012 \to 990$, $3774 \to 3514$) while also storing fewer. That
is a different effect from deferral, it is small, and the two searches compute different antichains
so the comparison is not like-for-like. It is not claimed as a result.

\paragraph{Consequence for the paper.} There is \textbf{one} algorithm here, not two.
\S\ref{sec:liveness} is the contribution; this section is the negative that says the
\texttt{explored} axis on \texttt{A[]} should not be funded further without a new idea about
\emph{selectivity}, which is the only door (\ref{eq:nodelay}) leaves open.
